% from ega3-7-fr.tex
\oldpage[III]{207}
il faut et il suffit que l'anneau $\Gamma(Z,\mathscr O_Z)$ ne soit pas un produit
de deux anneaux non réduits à 0.
\end{lemma}

\begin{proof}
Il est clair en effet que si $Z$ est réunion de deux ouverts non vides disjoints,
$\Gamma(Z,\mathscr O_Z)$ est isomorphe au produit des deux anneaux
$\Gamma(Z_1,\mathscr O_Z)$ et $\Gamma(Z_2,\mathscr O_Z)$ non réduits à 0.
Inversement, dire que $\Gamma(Z,\mathscr O_Z)$ est un tel produit équivaut à
dire qu'il y a dans $\Gamma(Z,\mathscr O_Z)$ un idempotent $s$ distinct de 0 et
de 1 ; pour tout $z\in Z$, $s_z$ est alors un idempotent dans $\mathscr O_z$,
donc égal à 0 ou 1. Mais il est clair que l'ensemble des $z$ tels que $s_z=0$
est ouvert ; d'autre part, si $s_z=1$, on a par définition $s(z)\ne0$, donc
l'ensemble des $z$ où $s_z=1$ est aussi ouvert ($0_{\mathrm I}$, 5.5.2) ; d'où
la conclusion.
\end{proof}

Il résulte de ce lemme que $X_y$ a exactement $m$ composantes connexes $X'_i$
et que $\Gamma(X'_i,\mathscr O_{X'_i})=\kappa(y)$ pour tout $i$. Comme $A$ a été
supposé local et artinien, son spectre est réduit à un point, donc $X$ et $X_y$
ont même espace sous-jacent ; $X$ a donc $m$ composantes connexes $X_i$ telles
que $X'_i=X_i\otimes_Y\kappa(y)$. On est ainsi finalement ramené au cas où
$R=\kappa(y)$ ; en vertu de (7.7.10, b)), on est ramené à prouver que
l'homomorphisme canonique
\[
  \Gamma(X,\mathscr O_X)
  \longrightarrow
  \Gamma(X_y,\mathscr O_{X_y})
\]
est surjectif ; mais cela est trivial, car le composé
\[
  \Gamma(Y,\mathscr O_Y)=A
  \longrightarrow
  \Gamma(X,\mathscr O_X)
  \longrightarrow
  \Gamma(X_y,\mathscr O_{X_y})=\kappa(y)
\]
est déjà surjectif.
\end{proof}

\begin{corollary}[7.8.7]
Sous les hypothèses de (7.8.6), il existe un voisinage ouvert $U$ de $y$ tel
que :
\begin{enumerate}
  \item[(i)] $f_*(\mathscr O_X)|U$ soit isomorphe à un
    $(\mathscr O_Y|U)$-Module de la forme $(\mathscr O_Y|U)^m$.
  \item[(ii)] Pour tout $z\in U$, l'homomorphisme canonique
    \[
      (f_*(\mathscr O_X))_z
      \otimes_{\mathscr O_z}\kappa(z)
      \longrightarrow
      \Gamma(X_z,\mathscr O_{X_z})
    \]
    est bijectif.
\end{enumerate}
\end{corollary}

\begin{proof}
(i) résulte de (7.8.6) et (7.8.4).

(ii) résulte de ce que $\mathscr T_0^U$ est exact (pour $U$ convenablement
choisi), et de (7.7.5.3).
\end{proof}

\begin{corollary}[7.8.8]
Supposons vérifiées les conditions de (7.8.6) et en outre que
$\Gamma(X_y,\mathscr O_{X_y})=\kappa(y)$. Alors il existe un voisinage ouvert
$U$ de $y$ tel que l'homomorphisme canonique
$\mathscr O_Y|U\to f_*(\mathscr O_X)|U$ soit bijectif.
\end{corollary}

\begin{proof}
En effet, il résulte de (7.8.7, (ii)) que l'entier $m$ figurant dans (7.8.7,
(i)) est nécessairement égal à 1.
\end{proof}

\begin{corollary}[7.8.9]
Sous les hypothèses de (7.8.6), il existe un voisinage ouvert $U$ de $y$, un
$\mathscr O_U$-Module cohérent $\mathscr Q$ (déterminé à isomorphisme unique
près) et un isomorphisme de foncteurs en le $\mathscr O_U$-Module quasi-cohérent
$\mathscr M$ :
\[
  \mathrm R^1f_*(f^*(\mathscr M))
  \xrightarrow{\sim}
  \mathscr{H}\!om_{\mathscr O_U}(\mathscr Q,\mathscr M).
  \tag{7.8.9.1}\label{III.7.8.9.1.fr}
\]
\end{corollary}

\begin{proof}
En effet, l'hypothèse entraîne que $\mathscr T_0^U$ est exact pour un $U$
convenable ; il suffit donc d'appliquer l'équivalence de (7.7.5, a)) et
(7.7.5, b$'$)) au cas $p=0$ et en prenant pour $\mathscr P_\bullet$ le complexe
réduit à son terme de degré 0 égal à $\mathscr O_X$.
\end{proof}

\begin{remarks}[7.8.10]
\textnormal{(i)} Sous les conditions de (7.8.6), considérons la factorisation de
Stein de $f$ (4.3.3)
\[
\xymatrix{
  X \ar[r]^{f'} & Y' \ar[r]^{g} & Y .
}
\]

